\section{The iterative procedure}
\label{sec:The-iterative-procedure}

The effect of backreaction is studied as an (infinite dimensional) fixed point problem, i.e. the solution $A_0$ to 
\begin{align}
    A_0 = f(A_0).
    \label{eq:fixed-point-problem}
\end{align}
Here, $f$ is a complicated update function, which for a given potential $A_0$
\begin{enumerate}
    \item calculates the modes of the scalar field $\phi$ from Eq.~\eqref{eq:mode-equation},
    \item calculates the corresponding vacuum polarization $\rho$ according to  Eq.~\eqref{eq:vacuum-polarization-long},
    \item returns the corresponding potential $A_0$ following Eq.~\eqref{eq:A0-general}.
\end{enumerate}
Self-consistent solutions, or fixed points of $f$, are looked for by studying the limit $\lim_{\kappa\to \infty}A_{0,\lambda}^{\kappa}$, where 
\begin{align}
    A_{0,\lambda}^{\kappa + 1} = f(A_{0,\lambda}^{\kappa}) =
    f^{\kappa}\left(A_{0, \lambda}^{0} \right), \textrm{with } f^\kappa =
    \underbrace{f \circ ... \circ f}_{k \textrm{ times}}.
    \label{eq:self-consistent-limit}
\end{align}
If the limit $A_{0, \lambda}^{\infty}$ exists, it is given the name
self-consistent potential, as well as the corresponding self-consistent vacuum
polarization $\rho_{\lambda}^{\infty}$, field modes $\phi_{n,
\lambda}^{\infty}$ and mode energies $\omega_{n, \lambda}^{\infty}$.

For low enough $\lambda\ll \lambda_c$, the initial value of the sequence $A_{0,
\lambda}^0$  can be safely set to the external field approximation $A_{0,
\lambda}^0(z) = -\lambda \left(z- \frac{1}{2}\right).$ For bigger
$\lambda$, this potential is no longer a valid starting potential  as it
might lead to non-convergence of $\lim_{\kappa \to \infty} A_0^{\kappa}$ or
complex $\omega_{n,\lambda}$. For higher values of $\lambda$,  $A_{0,
\lambda}^{0}$ is given by an already calculated self-consistent potential
$A_{0, \lambda'}^{\infty}$ with $\lambda = \lambda' + \Delta \lambda$, with
$\Delta \lambda > 0$. This scheme allows us to probe the $\lambda$-parameter
space.
 
For a given $\lambda$, at each step $\kappa$, the following steps are taken
\begin{subequations}
   \label{eq:iterative-procedure}
\begin{align}
    \frac{d {^2}\phi_{n, \lambda}^{\kappa}}{d {z^2}} {}  &= 
    -\left[ (\omega_{n, \lambda}^\kappa - \varepsilon A^{\kappa -1 }_{0, \lambda}(z) )^2 
     - a^2 m^2\right] \phi_{n, \lambda}^{\kappa}, 
    \label{eq:mode-equation-iterative}
    \\
    A^{\kappa }_{0, \lambda}(z) &= c A_{0, \lambda}^{\kappa-1} + \left(1-c
        \right)\left[ -\lambda\left(z - \frac{1}{2} \right) -
            \int_{\frac{1}{2}}^{z} \int_{0}^{z'} \rho_\lambda^{\kappa-1}(z'') d z'' d
    z' \right].
        \label{eq:A0-general-iterative}
\end{align}
\end{subequations}
Here, $c\in [0, 1]$ is introduced in order to ``relax" the procedure: For
$c=0$, the procedure is precisely the one used in~\cite[]{Ambjorn:1981xv}. For
$c$ close to 1, updates will only be corrections to the previous potential.
This avoids certain numerical instabilities as $\lambda   > \lambda_c$. Having
slower convergence schemes becomes increasingly relevant as $\lambda$ grows,
which makes the system more unstable.
%The mode subindices have been removed to avoid notation clutter.
%The iterative procedure is schematized in Fig.~\ref{fig:iterative-procedure},
%and it studies $\lim_{\kappa \to \infty} A_\lambda^\kappa $ for each $\lambda$.
%If the limit exists, it defines the self-consistent potential
%$A^{\infty}_\lambda$. Additionally, every other relevant quantity
%$\rho^\infty_\lambda, \omega_\lambda^{\infty}, \text{ and } \phi_\lambda
%^{\infty}$ also acquire the adjective ``self-consistent".

\begin{figure}[t]
\begin{center}
\begin{tikzpicture}[
    scale=0.87,
    box/.style = {draw, rectangle, minimum width=2cm, minimum height=1cm, fill={rgb:red,1;green,2;blue,2}, fill opacity=0.2, text opacity = 1},
    arrow/.style = {->, thick, >=Stealth},
    curved arrow/.style = {->, thick, >=Stealth, bend left}
]
% Nodes
\node[box] (A0) at (0, 0) {$A_{\lambda}^{0}$};
\node[box] (phi) at (3, 0) {$\omega_{n, \lambda}^\kappa, \phi_{n, \lambda}^\kappa$};
\node[box] (rho) at (8.25, 3) {$\rho_\lambda^\kappa$};
\node[box] (tildeA0) at (8.25, -3) {$A_\lambda^{\kappa}$};
\node (kappa) at (6.5, 0) {$\kappa \to \infty$};

% Arrows
\draw[arrow] (A0) -- (phi);

\draw[curved arrow, blue] (phi) to[out=50, in=135] (rho);
\draw[curved arrow, blue] (rho) to[out=55, in=135] (tildeA0);
\draw[curved arrow, blue] (tildeA0) to[out=45, in=130] (phi);

  %\draw[arrow] (phi) to[out=70, in=150, looseness=1] (rho);


% Central evolution arrow
\draw[->, thick, dashed, red] (8, 0) arc[start angle=360, end angle=0, radius=1.5cm];
%\draw[->, thick, dashed, red] (10, 0) arc[start angle=0, end angle=360, radius=3.5cm];

\end{tikzpicture}
\end{center}
\caption{Schematic of the iterative procedure.}%
\label{fig:iterative-procedure}
\end{figure}




%For a low $\lambda\ll \lambda_c$, the initial potential $A_\lambda^{0} = -\lambda
%\left(z-\frac{1}{2}\right)$ is given by the external field approximation. For
%higher values of $\lambda$, $A_\lambda^{0}$ is given by an already calculated
%self-consistent $A_{\lambda'}^{\infty}$, with $\lambda = \lambda' + \Delta
%\lambda$, for $\Delta \lambda > 0 $. 
